\documentclass{ctexart}

\usepackage{amsmath,amssymb,mathtools,amsthm,extarrows,color,enumitem,framed,amscd}
\usepackage[a4paper, left = 1cm, right = 1cm]{geometry}

\usepackage{tikz}
\usetikzlibrary{arrows.meta}
\usepackage{tikz-cd}

\usepackage[texMathDollars]{markdown}

\newcommand{\normal}{\mathrm{normal}}

% 取消首行缩进
\ctexset{autoindent=false}   % 先关闭 ctex 的自动调整功能
\setlength{\parindent}{0pt}  % 再把段首缩进设为 0

\usepackage{unicode-math}
%\setmathfont{STIX Two Math}
%\setmathfont{TeX Gyre Termes Math}
%\setmathfont{Libertinus Math}
%\setmathfont{Fira Math}

\begin{document}
\begin{markdown}
This conversation was generated with Duck.ai (https://duck.ai) using OpenAI's GPT-5.6 Luna Model. AI chats may display inaccurate or offensive information (see https://duckduckgo.com/duckai/privacy-terms for more info).

====================

User prompt 1 of 2 - 8/28/2026, 1:49:16 PM:
范畴论中的 equalizer, coequalizer 的定义是什么. 它们的直观是什么, 给出例子解释

GPT-5.6 Luna:
<!-- **Explaining Chinese definitions**

I need to provide an answer in Chinese while keeping it stable and not relying on the web. I should explain concepts with definitions, diagrams, intuition, and examples. Including the universal property sounds essential. I might discuss equalizer set functions and coequalizer quotients. It’s important to mention the idea of category duality, ensuring proper notation throughout. This way, I'll create a comprehensive and clear response! -->
在范畴论中，**equalizer（等化子）**和 **coequalizer（余等化子）**都是用来描述“使两个态射相等”的对象，但方式相反：

- equalizer：找出一个最大的子对象，使两个映射在其上相等；
- coequalizer：把源对象中的元素进行最小的商化，使两个映射在商对象上相等。

设有两个平行态射

$$
f,g:X\rightrightarrows Y.
$$

也就是 $f:X\to Y$ 和 $g:X\to Y$。

---

## 1. Equalizer：等化子

### 定义

如果有态射

$$
e:E\to X
$$

满足

$$
f\circ e=g\circ e,
$$

并且对任意态射 $h:Z\to X$，只要

$$
f\circ h=g\circ h,
$$

就存在唯一的态射

$$
u:Z\to E
$$

使得

$$
h=e\circ u,
$$

那么称 $e:E\to X$ 是 $f,g$ 的 **equalizer**。

图示为：
\end{markdown}
\[
\begin{tikzcd}
& E \arrow[d,"e"] \\
Z \arrow[ur,"u",dashed] \arrow[r,"h"] & X
\end{tikzcd}
\]
\begin{markdown}
其中要求存在唯一 $u:Z\to E$，使得 $h=e\circ u$。

更完整地说：
\end{markdown}
\[
\begin{tikzcd}
E \arrow[r,"e"] & X \arrow[r,shift left, "f"] \arrow[r,shift right,"g"'] & Y \\
Z \arrow[u,"u",dashed] \arrow[ur,"h"'] & &
\end{tikzcd}
\]
\begin{markdown}
并且 $E$ 满足

$$
f\circ e=g\circ e.
$$

### 直观理解

equalizer 是 $X$ 中那些“被 $f$ 和 $g$ 区分不出来”的部分。

在集合范畴 $\mathbf{Set}$ 中：

$$
E=\{x\in X\mid f(x)=g(x)\}.
$$

而 $e:E\hookrightarrow X$ 就是这个集合到 $X$ 的包含映射。

“泛”的意思是：任何落在 $X$ 中、并且使 $f$ 和 $g$ 相等的映射，都必须唯一地经过 $E$。

因此，equalizer 可以理解为：

> $X$ 中使 $f$ 与 $g$ 相等的最大部分。

这里的“最大”不是说元素数量最多，而是指具有泛性质，能够唯一接收所有其他满足条件的对象。

---

## 2. 集合中的例子

设

$$
X=\mathbb{R},\qquad Y=\mathbb{R},
$$

定义

$$
f(x)=x^2,\qquad g(x)=|x|.
$$

我们寻找满足

$$
x^2=|x|
$$

的实数 $x$。

解为：

$$
x=0,\quad x=1,\quad x=-1.
$$

所以 equalizer 是集合

$$
E=\{-1,0,1\},
$$

以及包含映射

$$
e:E\hookrightarrow \mathbb{R}.
$$

因为对于每个 $x\in E$，都有

$$
f(x)=g(x).
$$

如果另一个集合 $Z$ 有映射 $h:Z\to\mathbb{R}$，并且

$$
f\circ h=g\circ h,
$$

那么对每个 $z\in Z$，必有

$$
h(z)\in\{-1,0,1\}.
$$

因此 $h$ 唯一地分解为

$$
Z\xrightarrow{u}E\xrightarrow{e}\mathbb{R}.
$$

---

## 3. 群中的 equalizer

在群范畴 $\mathbf{Grp}$ 中，设有群同态

$$
f,g:G\to H.
$$

那么它们的 equalizer 是子群

$$
E=\{x\in G\mid f(x)=g(x)\}.
$$

特别地，如果 $g$ 是平凡映射，或者我们考虑

$$
f:G\to H,\qquad 0:G\to H,
$$

那么 equalizer 就是 $f$ 的核：

$$
\ker f=\{x\in G\mid f(x)=e_H\}.
$$

所以：

> 群同态的 kernel 是某个 equalizer 的特殊情形。

---

## 4. Coequalizer：余等化子

### 定义

如果有态射

$$
q:Y\to Q
$$

满足

$$
q\circ f=q\circ g,
$$

并且对任意态射 $h:Y\to Z$，只要

$$
h\circ f=h\circ g,
$$

就存在唯一的态射

$$
u:Q\to Z
$$

使得

$$
h=u\circ q,
$$

那么称 $q:Y\to Q$ 是 $f,g$ 的 **coequalizer**。

图示为：
\end{markdown}
\[
\begin{tikzcd}
X \arrow[r,shift left,"f"] \arrow[r,shift right,"g"'] & Y \arrow[r,"q"] \arrow[dr,"h"] & Q \arrow[d,"u",dashed]\\
& & Z
\end{tikzcd}
\]
\begin{markdown}
其中

$$
q\circ f=q\circ g,
$$

并且任何满足

$$
h\circ f=h\circ g
$$

的 $h:Y\to Z$，都唯一地经过 $q$。

### 直观理解

coequalizer 不再从 $X$ 中挑出一部分，而是把 $Y$ 中某些元素识别起来，使得 $f$ 和 $g$ 的结果变得相同。

在集合范畴中，coequalizer 大致是：

$$
Q=Y/{\sim},
$$

其中 $\sim$ 是由关系

$$
f(x)\sim g(x)\qquad (x\in X)
$$

生成的最小等价关系。

所以 coequalizer 可以理解为：

> 把 $f(x)$ 和 $g(x)$ 强制视为相同，从而得到的最一般商对象。

---

## 5. 集合中的 coequalizer 例子

设

$$
X=\{a,b\},\qquad Y=\{1,2,3,4\}.
$$

定义

$$
f(a)=1,\quad f(b)=2,
$$

以及

$$
g(a)=2,\quad g(b)=3.
$$

要求一个映射

$$
q:Y\to Q
$$

满足

$$
q\circ f=q\circ g.
$$

由 $a$ 得到：

$$
q(1)=q(2).
$$

由 $b$ 得到：

$$
q(2)=q(3).
$$

因此必须有

$$
q(1)=q(2)=q(3).
$$

元素 $4$ 不受任何关系约束，可以单独保留。

于是 coequalizer 可以取为

$$
Q=\{[1],[4]\},
$$

其中

$$
[1]=\{1,2,3\},\qquad [4]=\{4\}.
$$

商映射为

$$
q:Y\to Q,
$$

满足

$$
q(1)=q(2)=q(3)=[1],\qquad q(4)=[4].
$$

这就是把由 $f(x)$ 和 $g(x)$ 产生的元素强行合并。

---

## 6. 群中的 coequalizer

在群范畴中，设有两个群同态

$$
f,g:G\to H.
$$

那么 coequalizer 是一个商群

$$
q:H\to Q
$$

使得

$$
qf=qg.
$$

要求 $q(f(x))=q(g(x))$，等价于要求

$$
q\bigl(f(x)g(x)^{-1}\bigr)=e_Q.
$$

因此，应该把所有元素

$$
f(x)g(x)^{-1}
$$

都变成单位元，并且为了保持正规子群结构，需要取它们生成的正规闭包：

$$
N=\left\langle f(x)g(x)^{-1}\mid x\in G\right\rangle^{\normal}.
$$

于是 coequalizer 是自然商映射

$$
q:H\to H/N.
$$

特别地，如果 $g$ 是平凡同态，那么

$$
N=\langle f(G)\rangle^{\normal}.
$$

若 $f(G)$ 本身是正规子群，则 coequalizer 就是

$$
H/f(G).
$$

这说明：

> 在群范畴中，余等化子通常表现为“除以某个正规子群”。

---

## 7. Equalizer 与 coequalizer 的对偶性

它们是完全对偶的概念。

equalizer 的形式是：

$$
E\longrightarrow X\rightrightarrows Y,
$$

也就是先有一个对象 $E\to X$，使得两条路径相等。

coequalizer 的形式是：

$$
X\rightrightarrows Y\longrightarrow Q,
$$

也就是先有两条平行箭头，再从 $Y$ 出发把它们等化。

如果把范畴中的所有箭头反向：

- equalizer 会变成 coequalizer；
- coequalizer 会变成 equalizer。

因此：

$$\operatorname{Coeq}_{\mathcal C}(f,g) = \operatorname{Eq}_{\mathcal C^{op}}(f^{op},g^{op}).
$$

---

## 8. 与 kernel、cokernel 的关系

在有零对象的范畴中：

- kernel 是某种 equalizer；
- cokernel 是某种 coequalizer。

对于群或阿贝尔群：

$$
\ker(f)=\operatorname{Eq}(f,0),
$$

而在阿贝尔群范畴中：

$$
\operatorname{coker}(f)=\operatorname{Coeq}(f,0).
$$

具体地，对于群同态

$$
f:G\to H,
$$

kernel 是

$$
\ker f=\{x\in G\mid f(x)=e_H\},
$$

而 cokernel 是

$$
\operatorname{coker}(f)=H/\operatorname{im}(f)
$$

在阿贝尔群中成立；在一般群中需要取正规闭包。

---

## 9. 一句话总结

对于平行态射

$$
f,g:X\rightrightarrows Y,
$$

- **Equalizer**：

$$
E\to X
$$

是 $X$ 中使 $f$ 和 $g$ 相等的部分；在集合中就是

$$
\{x\in X\mid f(x)=g(x)\}.
$$

- **Coequalizer**：

$$
Y\to Q
$$

是把 $f(x)$ 和 $g(x)$ 识别为相同之后得到的商对象；在集合中就是

$$
Y/\bigl(f(x)\sim g(x)\bigr).
$$

可以记成：

> equalizer 是“取满足相等条件的子对象”，  
> coequalizer 是“施加相等关系得到商对象”。
\end{markdown}
\end{document}